\section{DISCUSSION AND CONCLUSIONS}
\label{sec:summary}

Finite-window emulation is limited by the history a model can access, not
by model capacity alone. Under our assumptions, sufficiently expressive
reset-window QELMs with exact expectations can approximate QRC outputs
arbitrarily well as their windows grow precisely when the influence of
the omitted past fades uniformly. We quantify the unavoidable error
caused by missing history and relate the context needed for a desired
accuracy to the rate of observable forgetting.

Conversely, our two-qubit construction preserves an old input's effect
on later outputs over arbitrarily long delays. Every fixed-window
predictor then faces a positive worst-case error floor that neither
additional features nor a longer finite window can remove. This
establishes a separation between carried memory and finite context,
rather than between quantum and classical computation. The guarantees
concern representability, not efficient learning or resource savings.

Analytical benchmarks and simulations distinguish three limitations:
missing history, insufficient features, and finite measurements. Their
practical message is to establish how much history the desired outputs
require before tuning how that history is represented and measured.
Missing information requires longer context or retained state;
increasing feature capacity alone cannot resolve it.

Together, these results provide a principled basis for designing and
comparing quantum temporal models. They identify when explicit context
can replace reservoir memory and when persistent observable information
prevents that replacement, making access to the past---not only the
richness of its representation---a central consideration in architecture
selection and resource allocation.
